% ============================================================
%  TEMA 3 — APARTADO 6: ALTERNATIVAS A LA REDUCCIÓN DE PLANTILLA
%                        E INCORPORACIÓN A LA EMPRESA
%
%  RECORDATORIO AL EQUIPO:
%  Antes de recurrir a despidos o EREs, las empresas pueden adoptar
%  medidas alternativas para equilibrar oferta y demanda de RRHH:
%
%  ALTERNATIVAS A LA REDUCCIÓN (cuando hay superávit):
%    1. Reducción de la jornada laboral (trabajo a tiempo parcial).
%    2. Reparto del trabajo (job sharing / worksharing).
%    3. Reducción salarial temporal (sacrificio compartido).
%    4. Congelación de contrataciones y de salarios.
%    5. Formación y recualificación (reskilling): reconvertir empleados
%       en perfiles demandados (p.ej. de cajero a asesor digital).
%    6. Movilidad funcional: asignar empleados a otras áreas con demanda.
%    7. Movilidad geográfica: trasladar empleados a zonas con necesidades.
%    8. Excedencias voluntarias y licencias sin sueldo.
%    9. Reducción de horas extra.
%   10. Externalización de funciones no estratégicas (outsourcing) para
%       liberar capacidad interna hacia tareas de mayor valor.
%
%  ALTERNATIVAS A LA INCORPORACIÓN (cuando hay déficit):
%  En lugar de contratar externamente, la empresa puede:
%    1. Promoción interna (movilidad vertical).
%    2. Formación acelerada para cubrir el gap de habilidades.
%    3. Trabajo en equipo / colaboración cross-funcional.
%    4. Horas extra (si el déficit es temporal y limitado).
%    5. Trabajo temporal o contratos de proyecto.
%    6. Externalización (BPO, outsourcing, contratas).
%    7. Automatización (RPA, IA) para reducir la necesidad de personal.
%    8. Acuerdos de colaboración con otras empresas del sector.
%    9. Internalización de funciones previamente externalizadas.
%   10. Programas de recontratación de exempleados (boomerang employees).
%
%  PARA ESTA SECCIÓN INDICAD:
%  - ¿Qué alternativas ha utilizado o utiliza CaixaBank?
%  - ¿Han sido efectivas? Ventajas e inconvenientes concretos.
%  - ¿Qué alternativas serían interesantes para CaixaBank y no usa aún?
% ============================================================

\section{Alternativas a la reducción de plantilla e incorporación a la empresa}

\subsection{Alternativas a la reducción de plantilla}

Antes de recurrir a medidas de ruptura laboral, CaixaBank puede —y debe—
explorar alternativas menos traumáticas que permitan ajustar la plantilla
preservando el empleo y el clima organizacional.

\begin{table}[H]
  \centering
  \small
  \caption{Alternativas a la reducción de plantilla en CaixaBank}
  \begin{tabularx}{\textwidth}{p{0.35\textwidth} p{0.12\textwidth} X}
    \toprule
    \textbf{Alternativa}                        & \textbf{¿Se aplica?} & \textbf{Observaciones} \\
    \midrule
    Reducción de jornada / trabajo parcial       & [Sí/No] & \\
    Congelación de contrataciones                & [Sí/No] & \\
    Reskilling / reconversión de perfiles        & [Sí/No] & \\
    Movilidad funcional interna                  & [Sí/No] & \\
    Movilidad geográfica                         & [Sí/No] & \\
    Excedencias voluntarias incentivadas         & [Sí/No] & \\
    Outplacement previo al despido               & [Sí/No] & \\
    Reducción de horas extra / subcontratación   & [Sí/No] & \\
    \bottomrule
  \end{tabularx}
\end{table}

\subsubsection{Reskilling y recualificación profesional}

% TODO: Este es probablemente el aspecto más relevante para CaixaBank.
% La digitalización bancaria hace que muchos perfiles de red de oficinas
% queden obsoletos. ¿Tiene CaixaBank un programa de reskilling para
% reconvertir a estos empleados en perfiles digitales?
% Mencionar si existe un programa concreto (p.ej. "Future Skills" o similar).

[Descripción a completar]

\textbf{Ventajas:}
\begin{itemize}
  \item Preserva el empleo y el conocimiento del negocio.
  \item Mejor imagen como empresa responsable.
  \item Menor coste que despedir y contratar externamente.
\end{itemize}

\textbf{Inconvenientes:}
\begin{itemize}
  \item No todos los empleados tienen el potencial o la motivación para
        reconvertirse en perfiles muy técnicos.
  \item Requiere inversión elevada en formación.
  \item El proceso es lento; no resuelve necesidades inmediatas de
        perfiles digitales especializados.
\end{itemize}

\subsubsection{Movilidad interna (funcional y geográfica)}

% TODO: ¿Tiene CaixaBank un sistema formal de movilidad interna?
% ¿Los empleados pueden postularse a vacantes internas? ¿Hay un portal
% de oportunidades internas en el SIRH? ¿Con qué frecuencia se producen
% traslados geográficos voluntarios o forzosos?

[Descripción a completar]

\subsection{Alternativas a la incorporación de nueva plantilla}

Cuando existe déficit de personal, CaixaBank no siempre recurre de inmediato
a la contratación externa. Las alternativas disponibles incluyen:

\subsubsection{Promoción interna}

% TODO: ¿Cuál es la política de CaixaBank respecto a la promoción interna?
% ¿Se publican primero las vacantes internamente antes de salir al mercado?
% ¿Tienen planes de carrera definidos?

[Descripción a completar]

\subsubsection{Automatización y digitalización}

% TODO: La RPA (Automatización Robótica de Procesos) y la IA están
% transformando el sector bancario. CaixaBank ha invertido fuertemente en
% tecnología. ¿Hay procesos que antes requerían empleados y ahora están
% automatizados? ¿Esto ha reducido la necesidad de incorporaciones en
% ciertas áreas?

[Descripción a completar]

\subsubsection{Externalización (\textit{outsourcing})}

% TODO: ¿Qué funciones tiene CaixaBank externalizadas?
% ¿IT, call centers, servicios de back-office, seguridad, limpieza?
% ¿Han seguido la tendencia de insourcing (recuperar funciones) en algún área?

[Descripción a completar]

\subsubsection{Trabajo temporal y contratos de proyecto}

[Descripción a completar]

\begin{notaequipo}
  Preguntad en la entrevista:
  \begin{enumerate}
    \item ¿Tienen algún programa de reskilling o upskilling para reconvertir
          perfiles tradicionales? ¿En qué consiste? ¿Cuántos empleados
          se han beneficiado?
    \item ¿Existe un portal de vacantes internas para la movilidad interna?
    \item ¿Cómo deciden entre cubrir una vacante internamente o acudir
          al mercado externo?
    \item ¿Qué procesos han automatizado recientemente con RPA o IA?
          ¿Ha supuesto una reducción de la plantilla en esas áreas?
    \item ¿Qué funciones tienen externalizadas? ¿Han internalizado alguna
          función en los últimos años?
  \end{enumerate}
\end{notaequipo}
